% Generated draft: requires mathematical and source review.
\documentclass[11pt]{article}
\usepackage[margin=25mm]{geometry}
\usepackage{fontspec}
\setmainfont{Latin Modern Roman}
\usepackage{amsmath,amssymb,mathtools,unicode-math}
\setmathfont{Latin Modern Math}
\usepackage{enumitem,needspace,xurl,hyperref}
\setlength{\parindent}{0pt}
\setlength{\parskip}{5pt}
\setlength{\emergencystretch}{3em}
\newcommand{\N}{\mathbb{N}}
\newcommand{\Z}{\mathbb{Z}}
\newcommand{\Q}{\mathbb{Q}}
\newcommand{\R}{\mathbb{R}}
\newcommand{\C}{\mathbb{C}}
\newcommand{\F}{\mathbb{F}}
\newcommand{\A}{\mathbb{A}}
\newcommand{\PP}{\mathbb{P}}
\newcommand{\E}{\mathbb{E}}
\newcommand{\eps}{\varepsilon}
\newcommand{\dd}{\,\mathrm{d}}
\begin{document}
\section*{UnsolvedMath ID: 801028}
\textbf{Source problem code:} catalogue entry
\subsection{Definitions and mathematical statement}
Let $S_\infty$ be the group of all permutations of $\N$, equipped with the topology of pointwise convergence, and let $X=\operatorname{Sub}(S_\infty)$ be the space of its closed subgroups, with the standard Borel structure inherited from the Effros Borel space of closed subsets of $S_\infty$. The Effros Borel structure is generated by the sets $\{F:F\cap U\ne\varnothing\}$ for open sets $U\subseteq S_\infty$.

For a precise conditional formulation, interpret group isomorphism as topological group isomorphism: write $G\cong_{\mathrm{top}}H$ if there is a bijective group homomorphism $G\to H$ that is continuous with continuous inverse. For $C\subseteq X$, let $E_C$ be this equivalence relation restricted to $C$. Under this interpretation, call $C$ invariant if $G\in C$, $H\in X$, and $G\cong_{\mathrm{top}}H$ imply $H\in C$. The supplied record does not specify whether isomorphism means topological or abstract group isomorphism, or whether invariant means closure under isomorphism or only under conjugation by elements of $S_\infty$; these choices remain qualifications on the formulation.

Let $\mathcal G=2^{[\N]^2}$ be the standard Borel space of simple undirected graphs with vertex set $\N$, and let $\cong_{\mathrm{gr}}$ denote graph isomorphism. For equivalence relations $E$ and $F$ on standard Borel spaces $Y$ and $Z$, respectively, write $E\leq_B F$ if there is a Borel map $f:Y\to Z$ such that
\[
\forall y,y'\in Y\quad \bigl(y\mathrel E y'\iff f(y)\mathrel F f(y')\bigr).
\]

With the preceding conditional conventions, the source asks:
\begin{enumerate}
\item Does there exist a Borel invariant class $C\subseteq X$ such that $E_C\not\leq_B\cong_{\mathrm{gr}}$?
\item Can the isomorphism relation for such a class be analytic complete?
\item What are the answers to the corresponding complexity questions for the class $C_{\mathrm{pro}}\subseteq X$ of pro-countable groups?
\end{enumerate}
Here a conventional topological meaning of pro-countable is being topologically isomorphic to an inverse limit
\[
\varprojlim_{n\in\N}D_n
=\left\{(d_n)_{n\in\N}\in\prod_{n\in\N}D_n:
\pi_n(d_{n+1})=d_n\text{ for every }n\in\N\right\},
\]
where each $D_n$ is a countable discrete group, each $\pi_n:D_{n+1}\to D_n$ is a group homomorphism, and the inverse limit has the subspace topology from the product. The source does not supply its definition of pro-countable, the intended indexing convention, or any requirement that bonding homomorphisms be surjective. It also does not explicitly assert that $C_{\mathrm{pro}}$ is Borel in the indicated coding, so that hypothesis must be checked before treating it as the Borel class requested in the first question.

The phrase analytic complete is also unresolved. One possible meaning is that $E_C\subseteq C\times C$ is analytic and every analytic subset $A$ of any standard Borel space $Y$ admits a Borel map $r:Y\to C\times C$ with $y\in A\iff r(y)\in E_C$. Another is completeness among analytic equivalence relations: $E_C$ is analytic and $F\leq_B E_C$ for every analytic equivalence relation $F$ on a standard Borel space. These meanings must not be silently identified; the supplied record does not select one.
\subsection{Short English statement}
The problem asks whether isomorphism on some Borel invariant class of closed permutation groups is too complex to admit a Borel reduction to countable graph isomorphism, whether analytic completeness can occur, and what happens for pro-countable groups. The meanings of isomorphism, invariance, analytic completeness, and pro-countability require confirmation from the source. The supplied registry labels the problem open.
\subsection{Sources}
\begin{itemize}
\item Local UnsolvedMath registry, permanent ID \(801028\), exact source problem code \texttt{catalogue entry}, title \emph{Open Problems in Mathematical Logic --- Problem 40}. The supplied record attributes the problem to Barmpalias, Gao et al., with proposed year 2026 and status \texttt{open}; these are supplied registry assertions.
\item Supplied source list: Barmpalias, Gao et al., \emph{Open Problems in Mathematical Logic} (2026), Problem 40, \url{https://arxiv.org/abs/2608.26628}. The background records access on 2026-10-16 and extraction method \texttt{arxiv\_latex}. This is supplied provenance, not a claim of independent source access or verification.
\item Supplied upstream revision: \texttt{d371ee917ac9647fbb06cd6fc529deec6f3dda82}. No separate literature assessment or \texttt{literature\_checked\_at} date is supplied.
\end{itemize}
\end{document}
